\scanpage{186}\mspage{42}
% archive sheet number: 186

$I'$, posant $\alpha=u(\alpha')$, la flèche induite $I'_{\alpha'}\longrightarrow I_\alpha$ soit une $W_0$-équivalence universelle, \deleted{\unsure{i.e. soit}} (cf critère a) pour cette notion).

c) Pour que $u$ soit une $W_0$-fibration faible, il \uline{faut} que $\forall i\in I$, l'inclusion $\underset{\unsure{\hookrightarrow|u}}{I'}_i\hookrightarrow{}_{i\backslash}I'$ soit une $W_0$-équivalence i.e. $\pi_0(I'_i)\unsure{\simeq|\stackrel{\sim}{\longrightarrow}}\pi_0({}_{i\backslash}I')$ soit bijectif, \uline{et} que $i\longmapsto\pi_0({}_{i\backslash}I')$ soit un préfaisceau loc. constant sur $I$ (i.e. $\underline{\pi}_0(I'/I)$ soit un \uline{revêtement} étale de $I$). Cette condition est aussi suffisante, si $u$ est $W_0$-lisse. \deleted{\unsure{ici}} De plus (pour mémoire) $u$ est une $W_0$-fibration faible ssi $u^{\mathrm{opp}}$ l'est.

La plus grande différence entre la notion de $W_0$-fibration, et celle de $W_0$-fibration faible, me paraît que les morphismes $W_0$-parfaits (i.e. $W_0$-lisses et $W_0$-propres), et a fortiori les $W$-parfaits \deleted{$W$-propres} pour un local. fond. $W\subset W_0$, sont forcément des $W$-fibrations faibles, mais (\uline{\unsure{$=$|même} si $W=W_\infty$}, le plus fin de tous les loc. fondamentaux) \unsure{ce|ne} ne sont pas des $\unsure{W|W_0}$-fibrations. L'exemple le plus frappant à cet égard est, sans doute celui des \unsure{revêtements|revêtement} étales, qui sont $W_\infty$-parfaits, (donc
